\section{Failure mode II: FISTA accelerates iterations but can destroy locality}
\label{sec:signed-star-fista-failure}

The previous failure was caused by too much order: residual nonnegativity
prevents accelerated motion.  Standard FISTA exhibits the opposite problem.
Its global iteration rate is accelerated, but its extrapolated trajectory can
touch a vertex whose degree is arbitrarily larger than the optimal support.

This section concerns the regularized RPPR objective
\begin{equation}
  F_\rho(x)
  =\frac12x^TQx-\alpha\langle D^{-1/2}s,x\rangle
   +\alpha\rho\|D^{1/2}x\|_1,
  \label{eq:signed-star-rppr-objective}
\end{equation}
and objective-gap target
$F_\rho(x_N)-F_\rho(x^\star)\le\varepsilon_{\rm obj}$.  It is therefore a
separate accuracy namespace from the semantic PPR theorem above.  The result
is the source-proved leaf-star obstruction of
\citet{fountoulakis2026complexity}, restated here with its decisive algebra.

\subsection{Standard strongly-convex FISTA}

For the source scaling, $Q$ has smoothness $L=1$ and
\eqref{eq:signed-star-rppr-objective} is $\alpha$-strongly convex.  Standard
fixed-momentum FISTA starts with $x_{-1}=x_0=0$ and uses
\begin{equation}
  \beta:=\frac{1-\sqrt\alpha}{1+\sqrt\alpha},
  \qquad
  y_k=x_k+\beta(x_k-x_{k-1}),
  \qquad
  x_{k+1}=\operatorname{prox}_{g_\rho}
      \bigl(y_k-\nabla f(y_k)\bigr).
  \label{eq:signed-star-fista}
\end{equation}
The proximal map soft-thresholds coordinate $i$ at
$\alpha\rho\sqrt{d_i}$.  The degree-weighted work model charges
\begin{equation}
  \operatorname{work}_k
  :=\operatorname{vol}(\operatorname{supp}y_k)
   +\operatorname{vol}(\operatorname{supp}x_{k+1}).
  \label{eq:signed-star-fista-work}
\end{equation}
Thus activating a single degree-$m$ coordinate costs at least $m$, even when
the number of active coordinates is only one.

\subsection{Leaf-seeded star with a one-vertex optimum}

Let $G(m)$ be a star with center $w$ of degree $m\ge2$, seed leaf $v$, and
$m-1$ other leaves.  Set
\begin{equation}
  \rho_0
  :=\frac{1-\alpha}{m(1+\alpha)+(1-\alpha)}.
  \label{eq:signed-star-fista-rho0}
\end{equation}
At this breakpoint, the unique RPPR optimizer is
\begin{equation}
  x^\star=x_v^\star e_v,
  \qquad
  x_v^\star=\frac{2\alpha(1-\rho_0)}{1+\alpha},
  \qquad
  \operatorname{supp}(x^\star)=\{v\}.
  \label{eq:signed-star-fista-optimum}
\end{equation}
To check this, the active KKT equation at $v$ yields the displayed value.
At the center,
\begin{equation}
  |\nabla_w f(x^\star)|
  =\frac{\alpha(1-\alpha)(1-\rho_0)}
          {(1+\alpha)\sqrt m}
  =\alpha\rho_0\sqrt m,
  \label{eq:signed-star-fista-tight-margin}
\end{equation}
so the center is inactive but has zero complementarity margin.  All other
leaves have zero smooth gradient.  Strong convexity makes the KKT point
unique.  Notice the geometry:
\begin{equation}
  \operatorname{vol}(\operatorname{supp}x^\star)=1,
  \qquad
  \operatorname{vol}(\partial\operatorname{supp}x^\star)=m.
  \label{eq:signed-star-fista-core-boundary}
\end{equation}
The optimizer is maximally local, while its one-vertex boundary is maximally
expensive.

\subsection{The exact overshoot}

For any extrapolated point supported on $\{v,w\}$, the center coordinate of
the forward point $u(y):=y-\nabla f(y)$ is
\begin{equation}
  u(y)_w
  =\frac{1-\alpha}{2}
      \left(y_w+\frac{y_v}{\sqrt m}\right).
  \label{eq:signed-star-fista-forward-center}
\end{equation}
Using the tight breakpoint identity
\eqref{eq:signed-star-fista-tight-margin}, the center becomes positive after
soft thresholding exactly when
\begin{equation}
  y_w\sqrt m+y_v>x_v^\star
  \label{eq:signed-star-fista-activation-criterion}
\end{equation}
for nonnegative $y_w,y_v$.  Before the center activates, $y_w=0$, so a seed
overshoot $y_v>x_v^\star$ is enough.

The preactivation dynamics are one-dimensional.  Put
\begin{equation}
  a:=\frac{1-\alpha}{2},
  \qquad
  e_k:=x_{k,v}-x_v^\star,
  \qquad
  \widetilde e_k:=y_{k,v}-x_v^\star.
  \label{eq:signed-star-fista-errors}
\end{equation}
As long as $w$ remains inactive, the soft threshold at the seed acts
affinely, giving
\begin{equation}
  e_{k+1}=a\bigl((1+\beta)e_k-\beta e_{k-1}\bigr).
  \label{eq:signed-star-fista-scalar-recurrence}
\end{equation}
The first errors are $e_0=-x_v^\star$ and $e_1=ae_0$.  Since
\begin{equation}
  a(1+\beta)=1-\sqrt\alpha,
  \qquad
  a(1+\beta)-\beta=\sqrt\alpha\,\beta,
  \label{eq:signed-star-fista-parameter-identities}
\end{equation}
we obtain
\begin{equation}
  \widetilde e_1=\sqrt\alpha\,\beta e_0<0,
  \qquad
  e_2=a\sqrt\alpha\,\beta e_0.
  \label{eq:signed-star-fista-first-errors}
\end{equation}
The second extrapolation reverses the sign:
\begin{align}
  \widetilde e_2
  &=(1+\beta)e_2-\beta e_1 \notag\\
  &=a\beta e_0\bigl((1+\beta)\sqrt\alpha-1\bigr)
    =-a\beta^2e_0
    =a\beta^2x_v^\star>0.
  \label{eq:signed-star-fista-exact-overshoot}
\end{align}
Thus $y_{2,v}>x_v^\star$, and
\eqref{eq:signed-star-fista-activation-criterion} forces
$x_{3,w}>0$.  Momentum has activated the degree-$m$ center after two
extrapolated steps, even though that coordinate is zero at the optimum.

\subsection{Work separation from ISTA}

The source turns this activation into an objective-gap lower bound.  Define
\begin{equation}
  \varepsilon_0(\alpha)
  :=\frac{\alpha^3(1-\alpha)^4\beta^4}
          {2(3+\alpha)^2}>0.
  \label{eq:signed-star-fista-epsilon0}
\end{equation}
For every
$0<\varepsilon_{\rm obj}\le\varepsilon_0(\alpha)$, the first four FISTA
iterates have not yet reached the target.  The verification uses the exact
one-dimensional gaps before activation and strong convexity afterward:
\begin{equation}
  F_{\rho_0}(x_ke_v)-F_{\rho_0}(x^\star)
  =\frac{1+\alpha}{4}e_k^2
  \quad(k=0,1,2),
  \label{eq:signed-star-fista-preactivation-gap}
\end{equation}
while $x_{3,w}>0$ and
$x_{3,v}-x_v^\star=a^2\beta^2x_v^\star$ imply a strictly positive
strong-convexity lower bound at $k=3$.  The choice
\eqref{eq:signed-star-fista-epsilon0} lies below all four quantities.

Iteration $k=2$ charges the support of $x_3$, which contains $w$; iteration
$k=3$ charges the support of
$y_3=x_3+\beta(x_3-x_2)$, which also contains $w$.  Consequently,
\begin{equation}
  W_{\rm FISTA}\ge2m.
  \label{eq:signed-star-fista-work-lower}
\end{equation}
By the source's imported ISTA support-containment result, ISTA remains inside
$\operatorname{supp}(x^\star)=\{v\}$ and reaches the same objective gap in
\begin{equation}
  W_{\rm ISTA}
  =O\!\left(
      \frac1\alpha\log\frac1{\varepsilon_{\rm obj}}
    \right),
  \label{eq:signed-star-ista-work-upper}
\end{equation}
independently of $m$.  Holding $\alpha$ and the permitted objective tolerance
fixed while sending $m\to\infty$ makes the FISTA/ISTA work ratio unbounded.

\subsection{What this failure does and does not say}

The obstruction does \emph{not} contradict FISTA's accelerated global
function-gap rate.  It says that iteration count does not control local work
when support volume varies along the trajectory.  Nor does it prove that all
momentum or all local acceleration fails: signed SOR on the center-star has
just succeeded, and the source also proves a conditional FISTA upper bound
when every spurious activation is confined to a controlled boundary set.

The exact lesson is narrower and more useful:
\begin{quote}
Small optimal support does not bound an accelerated method's transient
degree volume.  Any local FISTA theorem needs either confinement, a summable
spurious-work charge, or a safeguard that prevents expensive overshoot.
\end{quote}
